% généré par scripts/06_textes.py
\dosmeta{Année de classement : 341 (directe : nom consulaire lu). Références : ICUR I 1420 = ILCV 2816 (em) = ICVaticano p 299 = ICaRoma 89 = AE 1997, 166. Rattachement selon l’EDB : \emph{Incertae originis} (le volume I des ICVR réunit les inscriptions dont la provenance est perdue). Conservation selon l’EDB : \emph{Musei Vaticani: Lapidario Cristiano}. Support : Plaque de marbre. 49 × 186 cm (h. × l.). Inscription gravée. Signes décrits par l’EDB : chrisme. Reproductions signalées par l’EDB : dessin, ICVR I (de Rossi) et reproduction, Perret 1851 (droits réservés ; non reproduites). Fiche EDB : Filippo Piazzolla, première version 2015. Notice du catalogue p.~\pageref{not:4}.}
\txtbloc{Texte transmis (EDCS, état 2022)}{\lnum{1}Depossio Iuniani pri(die) Idus Apriles Marcellino et Probino conss(ulibus)\par
\lnum{2}qui ⸢v⸣ixit annis XL in pace dec⸢e⸣ssit et amator paupe\{ro\}rum ⸢v⸣ixit\par
\lnum{3}cum ⸢v⸣(i)rg(i)nia annis XV bene merenti ⸢v⸣irginia sua Bictora\par
\lnum{4}bene merenti fecit amatrix paupe\{ro\}rum et operaria}
\txtbloc{Texte de l’EDB (EDB31184)}{\lnum{1}depossio Iuniani pri(die) idus apriles Marcellino et Probino cons(ulibu)s\par
\lnum{2}qui bixit annis XL in pace dec⌜e⌝ssit et amator pauperorum bixit\par
\lnum{3}cum b⟨i⟩rg⟨i⟩nia annis XV benemerenti birginia sua Bictora\par
\lnum{4}benemerenti fecit amatrix pauperorum et operaria\par
\lnum{5}((monogramma Christi))}
